\section{Introduction}
\label{sec:intro}

Game designers seeking to create something new face a navigation problem: the design
space is vast, but the region they can perceive from their current position is limited
by the games they already know. Computational methods can extend that perceptual range
by mapping the space of existing games and identifying regions that are sparsely
populated --- not because those regions produce bad games, but because no one has
yet explored them.

\subsection{Design Space as a Navigable Geometry}

The TIGER framework~\cite{dellaLibera2026framework} places 150 games in a
five-dimensional space (Tension, Interaction, Game-Architecture, Elegance, Range).
This space is not uniformly populated: games cluster around a small number of
``genre attractors'' (high-complexity euros, light family games, social deduction
games), leaving large regions sparsely populated or empty. These sparse regions
are potential design opportunities --- places where a game could exist and might be
well-received, but where no game currently sits.

\subsection{From Evaluation to Generation}

Most evaluation frameworks are retrospective: they characterize games that already
exist. This paper demonstrates that TIGER can be used prospectively: by identifying
sparse regions in TIGER space and presenting them to experienced designers as explicit
design briefs, we test whether the computational analysis surfaces opportunities that
practitioners recognize as genuine and commercially viable.
